% GENERATED FILE. Edit canonical lessons, then run npm run generate:books.
% canonical-source-hash: fnv1a64:c61db2ed668de2d9
% canonical-lessons: BN-C521-bet, BN-C521-uiket, BN-C521-ran, BN-C521-kech, BN-C521-petano

\chapter{Bat and Ball}
\label{ch:bn-bat-and-ball}
\hlchaptermodality{521}

\begin{chapteropening}
\textbf{By the end of this chapter:} \emph{I can talk about bats, wickets, runs and catches.}

Complete the last lesson of chapter 521: I can talk about bats, wickets, runs and catches.
\end{chapteropening}

\section[\texorpdfstring{\bn{ব্যাট} (bæṭ)}{ব্যাট (bæṭ)}]{\bn{ব্যাট} (bæṭ) — a cricket bat}

\label{lesson:BN-C521-bet}



\begin{quote}
Before the new one: say the Bengali for a contestant, then the Bengali for to lead to victory.
\end{quote}

\subsection*{You'll want to know: \bn{ব্যাট}}
\textbf{\bn{ব্যাট}} — \emph{bæṭ} — \textquotedblleft{}a cricket bat\textquotedblright{}.

\begin{quote}
\bn{আমার} \bn{নতুন} \bn{ব্যাটটা} \bn{ভারী।}
\end{quote}

\begin{quote}
\emph{Āmār nôtun bæṭṭā bhārī.}
\end{quote}

\begin{quote}
My new bat is heavy.
\end{quote}

\begin{cousinweb}
\textbf{\bn{ব্যাট}} (\emph{bæṭ}) is borrowed from English \emph{bat}.
\end{cousinweb}

\begin{grammarlens}[title={What you've built}]
One more word for describing what you lived and felt.
\end{grammarlens}

\subsection*{Your turn}
\begin{itemize}
\raggedright
  \item \emph{Say it:} \emph{bæṭ}
  \item \emph{Say it:} \emph{bæṭ}, once more
  \item \emph{Say it:} \emph{Āmār nôtun bæṭṭā bhārī.} — My new bat is heavy.
  \item \emph{Recall:} say the Bengali for a contestant, then the Bengali for to lead to victory, then say \emph{bæṭ} again
\end{itemize}

\begin{tcolorbox}[breakable,colback=teal!4,colframe=teal!35!black,title={Before you move on}]
What does \bn{ব্যাট} mean? (\textquotedblleft{}A cricket bat\textquotedblright{}.) Say it once more.
\end{tcolorbox}
\section[\texorpdfstring{\bn{উইকেট} (uikeṭ)}{উইকেট (uikeṭ)}]{\bn{উইকেট} (uikeṭ) — a wicket}

\label{lesson:BN-C521-uiket}



\begin{quote}
Before the new one: say the Bengali for to lead to victory, then the Bengali for a cricket bat.
\end{quote}

\subsection*{You'll want to know: \bn{উইকেট}}
\textbf{\bn{উইকেট}} — \emph{uikeṭ} — \textquotedblleft{}a wicket\textquotedblright{}.

\begin{quote}
\bn{বলটা} \bn{উইকেটে} \bn{লাগল।}
\end{quote}

\begin{quote}
\emph{Bôlṭā uikeṭe lāglô.}
\end{quote}

\begin{quote}
The ball hit the wicket.
\end{quote}

\begin{cousinweb}
\textbf{\bn{উইকেট}} (\emph{uikeṭ}) is borrowed from English \emph{wicket}.
\end{cousinweb}

\begin{grammarlens}[title={What you've built}]
One more word for describing what you lived and felt.
\end{grammarlens}

\subsection*{Your turn}
\begin{itemize}
\raggedright
  \item \emph{Say it:} \emph{uikeṭ}
  \item \emph{Say it:} \emph{uikeṭ}, once more
  \item \emph{Say it:} \emph{Bôlṭā uikeṭe lāglô.} — The ball hit the wicket.
  \item \emph{Recall:} say the Bengali for to lead to victory, then the Bengali for a cricket bat, then say \emph{uikeṭ} again
\end{itemize}

\begin{tcolorbox}[breakable,colback=teal!4,colframe=teal!35!black,title={Before you move on}]
What does \bn{উইকেট} mean? (\textquotedblleft{}A wicket\textquotedblright{}.) Say it once more.
\end{tcolorbox}
\section[\texorpdfstring{\bn{রান} (rān)}{রান (rān)}]{\bn{রান} (rān) — a run in cricket}

\label{lesson:BN-C521-ran}



\begin{quote}
Before the new one: say the Bengali for a cricket bat, then the Bengali for a wicket.
\end{quote}

\subsection*{You'll want to know: \bn{রান}}
\textbf{\bn{রান}} — \emph{rān} — \textquotedblleft{}a run in cricket\textquotedblright{}.

\begin{quote}
\bn{জিততে} \bn{আমাদের} \bn{দশ} \bn{রান} \bn{দরকার।}
\end{quote}

\begin{quote}
\emph{Jitte āmāder dôsh rān dôrkār.}
\end{quote}

\begin{quote}
We need ten runs to win.
\end{quote}

\begin{cousinweb}
\textbf{\bn{রান}} (\emph{rān}) is borrowed from English \emph{run}, used only for the scoring run in cricket.
\end{cousinweb}

\begin{grammarlens}[title={What you've built}]
One more word for describing what you lived and felt.
\end{grammarlens}

\subsection*{Your turn}
\begin{itemize}
\raggedright
  \item \emph{Say it:} \emph{rān}
  \item \emph{Say it:} \emph{rān}, once more
  \item \emph{Say it:} \emph{Jitte āmāder dôsh rān dôrkār.} — We need ten runs to win.
  \item \emph{Recall:} say the Bengali for a cricket bat, then the Bengali for a wicket, then say \emph{rān} again
\end{itemize}

\begin{tcolorbox}[breakable,colback=teal!4,colframe=teal!35!black,title={Before you move on}]
What does \bn{রান} mean? (\textquotedblleft{}A run in cricket\textquotedblright{}.) Say it once more.
\end{tcolorbox}
\section[\texorpdfstring{\bn{ক্যাচ} (kæch)}{ক্যাচ (kæch)}]{\bn{ক্যাচ} (kæch) — a catch in cricket}

\label{lesson:BN-C521-kech}



\begin{quote}
Before the new one: say the Bengali for a wicket, then the Bengali for a run in cricket.
\end{quote}

\subsection*{You'll want to know: \bn{ক্যাচ}}
\textbf{\bn{ক্যাচ}} — \emph{kæch} — \textquotedblleft{}a catch in cricket\textquotedblright{}.

\begin{quote}
\bn{অধিনায়ক} \bn{একটা} \bn{দুর্দান্ত} \bn{ক্যাচ} \bn{ধরল।}
\end{quote}

\begin{quote}
\emph{Ôdhināyôk ekṭā durdāntô kæch dhôrlô.}
\end{quote}

\begin{quote}
The captain took a brilliant catch.
\end{quote}

\begin{cousinweb}
\textbf{\bn{ক্যাচ}} (\emph{kæch}) is borrowed from English \emph{catch}; with \bn{ধরা}: \bn{ক্যাচ} \bn{ধরা}, to take a catch.
\end{cousinweb}

\begin{grammarlens}[title={What you've built}]
One more word for describing what you lived and felt.
\end{grammarlens}

\subsection*{Your turn}
\begin{itemize}
\raggedright
  \item \emph{Say it:} \emph{kæch}
  \item \emph{Say it:} \emph{kæch}, once more
  \item \emph{Say it:} \emph{Ôdhināyôk ekṭā durdāntô kæch dhôrlô.} — The captain took a brilliant catch.
  \item \emph{Recall:} say the Bengali for a wicket, then the Bengali for a run in cricket, then say \emph{kæch} again
\end{itemize}

\begin{tcolorbox}[breakable,colback=teal!4,colframe=teal!35!black,title={Before you move on}]
What does \bn{ক্যাচ} mean? (\textquotedblleft{}A catch in cricket\textquotedblright{}.) Say it once more.
\end{tcolorbox}
\section[\texorpdfstring{\bn{পেটানো} (peṭāno)}{পেটানো (peṭāno)}]{\bn{পেটানো} (peṭāno) — to beat; to hit hard}

\label{lesson:BN-C521-petano}



\begin{quote}
Before the new one: say the Bengali for a run in cricket, then the Bengali for a catch in cricket.
\end{quote}

\subsection*{You'll want to know: \bn{পেটানো}}
\textbf{\bn{পেটানো}} — \emph{peṭāno} — \textquotedblleft{}to beat; to hit hard\textquotedblright{}.

\begin{quote}
\bn{অধিনায়ক} \bn{বলটা} \bn{মাঠের} \bn{বাইরে} \bn{পেটাল।}
\end{quote}

\begin{quote}
\emph{Ôdhināyôk bôlṭā māṭher bāire peṭālô.}
\end{quote}

\begin{quote}
The captain hit the ball right out of the ground.
\end{quote}

Past (he/she …): \bn{তিনি} \bn{বলটা} \bn{পেটালেন}, \emph{tini bôlṭā peṭālen}.

\begin{cousinweb}
\textbf{\bn{পেটানো}} (\emph{peṭāno}) is a tadbhava: CDIAL lists Bengali \emph{piṭā}, \textquotedblleft{}to beat\textquotedblright{}, under Sanskrit \emph{piṭṭayati}, \textquotedblleft{}stamps into a solid mass\textquotedblright{} (CDIAL 8165).
\end{cousinweb}

\begin{grammarlens}[title={What you've built}]
One more word for describing what you lived and felt.
\end{grammarlens}

\subsection*{Your turn}
\begin{itemize}
\raggedright
  \item \emph{Say it:} \emph{peṭāno}
  \item \emph{Say it:} \emph{peṭāno}, once more
  \item \emph{Say it:} \emph{Ôdhināyôk bôlṭā māṭher bāire peṭālô.} — The captain hit the ball right out of the ground.
  \item \emph{Recall:} say the Bengali for a run in cricket, then the Bengali for a catch in cricket, then say \emph{peṭāno} again
\end{itemize}

\begin{tcolorbox}[breakable,colback=teal!4,colframe=teal!35!black,title={Before you move on}]
What does \bn{পেটানো} mean? (\textquotedblleft{}To beat; to hit hard\textquotedblright{}.) Say it once more.
\end{tcolorbox}
